%% kapa: exam=1 round=22 year=2011 period=1 form=A subject=civil_law subject_ko=민법 seq=16 number=16
%% source: Q-Net 감정평가사 기출문제 (cst00309) artlSeq=5036790 file=extracted/1008/5036790/2011년 제22회 감정평가사 1차시험 기출문제/1교시.hwp
%% conversion: kordoc 4.24.51 -> tools/gen1.py (auto); review=proofread-sonnet; flags=none
\begin{question}{16}
물권적 청구권에 관한 설명으로 옳지 않은 것은? (다툼이 있으면 판례에 의함)
\begin{choices}
  \item 소유자는 자신의 물건을 권원 없이 점유하는 자에 대해 점유자가 과실이 없더라도 그 반환을 청구할 수 있다.
  \item 물권적 청구권은 물권과 분리하여 양도할 수 있다.
  \item 진정명의회복을 위한 소유권이전등기청구권은 소유물 방해배제청구권에 해당한다.
  \item 공사로 인하여 점유의 방해를 받은 경우, 공사착수 후 1년이 경과한 때에는 방해제거를 청구할 수 없다.
  \item 지역권자는 지역권 침해를 이유로 승역지의 반환을 청구할 수 없다.
\end{choices}
\end{question}
\answer{2}
